\section{Related Work}
\label{sec:related}

\subsection{Position Bias in LLMs}

Large language models exhibit systematic position bias when processing long contexts. \citet{Liu2023LostIT} identified the ``lost-in-the-middle'' phenomenon, demonstrating that LLMs perform best when relevant information appears at the beginning or end of the input, with significant performance degradation for middle positions. This U-shaped attention pattern has been attributed to intrinsic attention bias, where tokens at extreme positions receive disproportionately high attention regardless of their semantic relevance~\citep{Hsieh2024FoundIT}. \citet{Yi2025AttentionBW} characterized this as the ``attention basin'' phenomenon, showing that shallow attention layers play a critical role in determining which positions receive attention. Separately, \citet{Xiao2023EfficientSL} discovered ``attention sinks''---initial tokens that absorb excess attention mass even when semantically unimportant---which enables efficient streaming inference but also contributes to position-dependent attention allocation.

\subsection{Context Window Extension}

Extending the context window of pretrained LLMs has been an active research area. Rotary Position Embedding (RoPE)~\citep{Su2021RoFormerET} encodes absolute positions using rotation matrices while incorporating relative position information in self-attention, enabling flexible sequence lengths. However, RoPE-based models fail to generalize beyond their training sequence length. Position Interpolation~\citep{Chen2023ExtendingCW} addresses this by linearly down-scaling position indices to fit within the original context window, extending LLaMA models to 32K tokens with minimal fine-tuning. YaRN~\citep{Peng2023YaRNEC} improves upon this with a compute-efficient extension method requiring 10$\times$ fewer tokens and 2.5$\times$ fewer training steps, enabling context windows up to 128K tokens. More recently, LongRoPE~\citep{Ding2024LongRoPEEL} extends context windows to 2M tokens through progressive extension and non-uniform positional interpolation. These methods focus on enabling longer contexts but do not directly address the position bias that affects information utilization within those contexts.

\subsection{Attention-Based Document Ranking}

Several methods leverage attention patterns for document ranking in retrieval-augmented generation. Attention Sorting~\citep{Peysakhovich2023AttentionSC} iteratively reorders documents based on the attention they receive during generation, placing high-attention documents at the end where recency-biased models attend more effectively. \citet{Chen2024AttentionIL} propose in-context re-ranking (ICR), which uses attention pattern changes caused by the query for efficient zero-shot re-ranking with only two forward passes. To address inconsistencies in LLM-based ranking, \citet{Zeng2024LLMRankFusionMI} introduce LLM-RankFusion, which mitigates order and transitive inconsistencies through calibration and ranker aggregation. \citet{Tang2023FoundIT} propose permutation self-consistency, which marginalizes over different list orderings to produce order-independent rankings. Our work builds on Attention Sorting but investigates whether explicit position-bias correction can reduce the number of required sorting iterations.

